% F5d -- How one training run actually works, step by step.
\begin{figure}[!htbp]
\centering
\fitwidth{%
\begin{tikzpicture}[font=\small,
  st/.style={draw=clrule,fill=clsoft,rounded corners=3pt,align=center,
             inner sep=6pt,text width=2.85cm,minimum height=2.75cm},
  ar/.style={-{Stealth[length=5pt]},clrule,very thick},
  num/.style={circle,draw=clmodel,fill=white,inner sep=1.2pt,
              font=\scriptsize\bfseries,text=clmodel},
]
\node[st] (a) at (0,0)
  {{\footnotesize\bfseries\color{clmodel} Take a slice}\\[2pt]\footnotesize
   a stretch of one training person's day: hour, antenna, hour, antenna};
\node[st] (b) at (3.55,0)
  {{\footnotesize\bfseries\color{clmodel} Hide the next antenna}\\[2pt]\footnotesize
   the model reads the slice and names what it thinks comes next};
\node[st] (c) at (7.10,0)
  {{\footnotesize\bfseries\color{clmodel} Measure the surprise}\\[2pt]\footnotesize
   compare with the antenna really visited; a confident wrong answer costs more
   than a hesitant one};
\node[st] (d) at (10.65,0)
  {{\footnotesize\bfseries\color{clmodel} Nudge the small part}\\[2pt]\footnotesize
   only the added numbers move; the 500 million original ones never do};

\draw[ar] (a) -- (b); \draw[ar] (b) -- (c); \draw[ar] (c) -- (d);
\node[num] at (a.north west) {1}; \node[num] at (b.north west) {2};
\node[num] at (c.north west) {3}; \node[num] at (d.north west) {4};

% --- the loop back --------------------------------------------------------
\draw[ar] (d.south) -- ++(0,-0.65) -- ++(-10.65,0) -- (a.south);
\node[font=\scriptsize,text=clrule] at (5.32,-2.42)
  {repeat, slice after slice, until every training person has been read once:
   that is \textbf{one pass}};

% --- the checkpoint rule --------------------------------------------------
\node[draw=clgain,fill=clgain!10,rounded corners=3pt,align=center,inner sep=6pt,
      text width=13.4cm,font=\small,anchor=north] (ck) at (5.32,-2.85)
  {{\footnotesize\bfseries\color{clgain} 5. At the end of every pass: grade the
    model on people it is not training on, save a copy, and keep only the best
    copy of the whole run}\\[2pt]\footnotesize
   On this task the best copy arrives between the fourth and the ninth pass, and
   everything after it is worse. Keeping the last copy instead of the best one
   costs 15.5 correct answers per hundred (Chapter~\ref{ch:first}).};
\end{tikzpicture}}
\caption{How one training run works}
\label{fig:training-loop}
\figtext{%
\figpara{Steps 1 to 4 are the whole of training.}{A slice of somebody's day is
shown to the model with the next antenna hidden. The model names an antenna. The
answer is compared with the truth and turned into a single number measuring how
surprised the model was, so that being confidently wrong costs more than
hesitating. That number is then used to nudge the model's adjustable numbers
very slightly in the direction that would have made it less wrong. Then the next
slice.}

\figpara{Only a small part of the model ever moves.}{Step 4 is where the method
of Figure~\ref{fig:lora} bites. The 500 million numbers the model was published
with are frozen and never touched; the nudging applies only to a small set of
numbers added beside them, about 2 out of every 100. That is what makes a
training run fit in half an hour on a free graphics card.}

\figpara{Step 5 is not an optimisation, it is a safety rule.}{Left to itself, a
model trained on a few hundred short days stops learning the habits of the
population and starts learning the days themselves by heart. When that happens
its score on people it has never seen turns around and falls, while its score on
its own training data keeps improving, so the training figures alone will happily
report success for another forty passes. Grading on separate people after each
pass, and keeping the best copy rather than the last, is what prevents that, and
it was the largest single improvement of the whole internship.}}
\end{figure}
